\corrigesChapitre{7}{PGCD, Bézout et Gauss}
\begin{corrige}{1}{Calculer un PGCD}
\begin{questions}
\item $48=2\times18+12$, $18=12+6$, $12=2\times6$ : PGCD $6$.
\item $35=2\times12+11$, $12=11+1$ : PGCD $1$.
\item On prend les valeurs absolues : PGCD $6$.
\end{questions}
\end{corrige}
\begin{corrige}{2}{Vérifier Bézout}
\begin{questions}
\item $1008-990=18$.
\item Elle prouve que le PGCD divise $18$, mais il faut aussi vérifier que $18$ divise les deux nombres.
\item $252=14\times18$ et $198=11\times18$ ; $18$ est un diviseur commun et tout diviseur commun divise $18$, donc c’est le PGCD.
\end{questions}
\end{corrige}
\begin{corrige}{3}{Un inverse}
\begin{questions}
\item $15=2\times7+1$.
\item Multiplier par $5$ : $x\equiv20\equiv6\pmod7$.
\item $x=6$.
\end{questions}
\end{corrige}
\begin{corrige}{4}{Équation impossible}
\begin{questions}
\item $3$.
\item Non, car le membre de gauche est toujours divisible par $3$.
\item Oui : $x=-1,y=1$ donne $-6+9=3$.
\end{questions}
\end{corrige}
\begin{corrige}{5}{Remonter Euclide}
\begin{questions}
\item $35=2\times12+11$, $12=11+1$.
\item $1=12-(35-2\times12)=3\times12-35$.
\item $3$ est un inverse de $12$ modulo $35$.
\end{questions}
\end{corrige}
\begin{corrige}{6}{Gauss en action}
\begin{questions}
\item $\operatorname{PGCD}(7,15)=1$ ; Gauss s’applique.
\item Non : $n=3$ fournit un contre-exemple.
\item $6\mid4n$ implique $3\mid2n$ ; comme $3$ et $2$ sont premiers entre eux, $3\mid n$.
\end{questions}
\end{corrige}
\begin{corrige}{7}{Congruence avec PGCD non trivial}
\begin{questions}
\item Diviser par $3$ donne $2x\equiv1\pmod3$ ; donc $x\equiv2\pmod3$.
\item Les représentants sont $2,5,8$.
\item Le module réduit est $3$ ; trois classes modulo $9$ correspondent à la même classe modulo $3$.
\end{questions}
\end{corrige}
\begin{corrige}{8}{Toutes les solutions entières}
\begin{questions}
\item Diviser par $3$ : $5x+7y=2$.
\item $(x_0,y_0)=(-1,1)$.
\item $x=-1+7k$, $y=1-5k$, $k\in\Z$, d’après Gauss après soustraction de la solution particulière.
\end{questions}
\end{corrige}
\begin{corrige}{9}{Une contrainte positive}
\begin{questions}
\item Une solution est $(2,5)$ ; toutes sont $x=2+5k,y=5-3k$.
\item $k\geq0$ et $k\leq1$, donc $k=0$ ou $1$.
\item Deux solutions : $(2,5)$ et $(7,2)$. Les contraintes ont réduit une famille infinie à deux couples.
\end{questions}
\end{corrige}
\begin{corrige}{10}{Inverser modulo $26$}
\begin{questions}
\item $7\times15=105=4\times26+1$, donc inverse $15$.
\item $x\equiv75\equiv23\pmod{26}$.
\item $7\times23=161=6\times26+5$.
\end{questions}
\end{corrige}
\begin{corrige}{11}{Deux congruences, problème guidé}
\begin{questions}
\item $2+3k\equiv3$ donne $3k\equiv1\pmod5$, donc $k\equiv2\pmod5$.
\item $k=2+5t$, donc $x=8+15t$, $t\in\Z$.
\item Ajouter $15$ conserve les restes modulo $3$ et $5$ ; la différence de deux solutions est divisible par ces deux entiers premiers entre eux, donc par $15$.
\end{questions}
\end{corrige}
\begin{corrige}{12}{Prouver Gauss}
\begin{questions}
\item Il existe $u,v\in\Z$ tels que $au+bv=1$.
\item $auc+bvc=c$.
\item Les deux termes de gauche sont divisibles par $a$, donc $a\mid c$. Le caractère premier entre eux est l’hypothèse qui permet l’identité égale à $1$.
\end{questions}
\end{corrige}
